This chapter will focus on presenting the context of this thesis.
In particular, in \hyperref[sec:fuzz_background]{Section 1.1} the fundamentals of fuzzing will be discussed, with particular attention to coverage-guided grey-box fuzzing.
In \hyperref[sec:biniary_rewriting]{Section 1.2} the background on binary rewriting is discussed, along with its approaches and difficulties.
In the next chapther will be presented the approach used in order to resolve the main problems regarding the solution 
discussed in this thesis. In order to do so some concept of binary loading must be presented (\hyperref[sec:loading_relocations]{Section 1.3}).
In order to better understand a motivation to a design choice expressed in \cref{sec:actual_instrumentation}, in 
\hyperref[sec:preamble]{Section 1.4} are presented two of the compiler optimization in terms of function preamble.
\hyperref[sec:problem_statement]{Section 1.5} states the problem addressed by this work,
and \hyperref[sec:state_of_the_art]{Section 1.6} reviews the state of the art in binary-only fuzzing, also focusing on a description
of AFL/AFL++ internals. In the last section (\hyperref[sec:goals-and-challenges]{Section 1.7}) I will present the main
goals and challenges of this thesis, and explain why the implemented tool differs from the state of the art.

\section{Fuzzing Background} \label{sec:fuzz_background}
    Fuzzing is nowadays a widely used approach for software testing and vulnerability assessment \cite{Zhu2022,Mallissery2023}.
    Its main strength derives from the fact that it can rapidly produce various inputs to be fed into a target executable.
    In general, a fuzzer can be divided into the \textbf{Input Generator}, which is responsible for providing inputs to the \textbf{Executor} that executes them.
    During the execution, the \textbf{Defect Monitor} checks whether the executed input causes a crash of the target.\\
    In particular, the Input Generator must choose which input to provide to the executor among all the inputs in the \textit{Input Space}.
    A subset of the \textit{Input Space} is represented by the \textit{Valid Input Space}.
    This is a subset of the input space that specifically contains all the inputs that are well formed with respect to the
    application being tested.
    For example, if the fuzzed target is a PDF reader, a valid input is a PDF file.
    A good Input Generator must choose inputs that are also valid inputs, because only by doing this is there a probability of finding a vulnerability \cite{Zhu2022}.
    Hence the need for a precise input selection and mutation strategy arises, in order to discover as many edge cases as possible.\\
    In order to do so, the Input Generator must obtain some information from the Executor.
    Depending on the amount of information the Input Generator has, fuzzing can be classified into
    \textit{White-box}, \textit{Black-box} or \textit{Grey-box}.
    In Black-box fuzzing, the fuzzer does not have any knowledge of the execution \cite{Miller1990}.
    In White-box fuzzing, instead, the source code must be available in the general case, and techniques such as Symbolic and Concolic execution are used \cite{Godefroid2012}.
    In Grey-box fuzzing, instead, the inputs are chosen using, for example, \textbf{Code Coverage} information (i.e. the number of lines of code visited during execution of the target).
    This approach is called \textit{Coverage-Guided Grey-box fuzzing} and is based on prioritizing the inputs that discover new sections of code.
    Since the source code is not always available, Code Coverage is not always a viable Coverage Measure.
    In general, the coverage measures can be of different granularity \cite{CollAFL}.
    The three main options are:
    \begin{itemize}
        \item \textbf{Block Coverage}: Keeps track of the hit count for each \gls{BB}. This approach is used by some
            popular fuzzers such as libFuzzer or honggfuzz \cite{CollAFL}. This method is often too imprecise because
            it fails to capture the transitions between BBs.
        \item \textbf{Edge Coverage}: Keeps track of the hit count of every edge in the \gls{CFG} (i.e. the transition between two \gls{BB}s).
            This is the approach used by AFL and AFL++ \cite{aflplusplus}. It maintains overheads similar to Block Coverage but can
            track the transitions between \gls{BB}s and not only their execution, adding some context information.
            However, it lacks precision in keeping track of the order in which the edges are taken.
        \item \textbf{Path Coverage}: This is the most complete metric and keeps track of the complete path taken during the execution
            of a specific input, thus also keeping track of the order in which the edges are taken.
            The main drawback of this metric is that a single execution path could become very large, resulting in
            large runtime and memory overheads.
    \end{itemize}

    One of the best approach to Code Coverage, in terms of tradeoff between overhead and effectiveness, 
    is \textbf{Edge Coverage} \cite{CollAFL}, which keeps track of which edges
    in the \gls{CFG} are taken during execution without the overheads incurred by keeping track of the entire execution path.
    Coverage information is obtained through the instrumentation of the target application.
    Depending on the availability of the source code, the instrumentation can be added in different ways:
    \begin{itemize}
        \item The source code is available: In this case the instrumentation is added at \textit{Compile Time}.
        \item The source code is not available: In this case there are two different techniques for instrumenting the application:
        \begin{itemize}
            \item \textbf{\gls{SBI}}: The instrumentation is added before the target is executed, without access to the
                source code, relying on Static Binary Rewriting techniques.
            \item \textbf{\gls{DBI}}: The instrumentation is added during program execution, relying on Dynamic Binary Rewriting techniques.
        \end{itemize}
    \end{itemize}
    Regardless of the availability of the source code, the \textit{Coverage-Guided Grey-box fuzzing} technique can be summarized
    as shown in \cref{fig:coverageguidedfuzzing}.
    \begin{figure}[ht]
        \caption{Generic diagram describing the Coverage-Guided Grey-box fuzzing workflow}
        \centering
        \includegraphics[width=\textwidth]{CoverageGuidedGrayBoxFuzzing.png}
        \label{fig:coverageguidedfuzzing}
    \end{figure}
    \\Specifically for \textbf{Edge Coverage}, \cref{fig:coverageguidedfuzzing} can be described as follows:
    \begin{enumerate}
        \item \textbf{Seed Corpus}: Represents the set of inputs on which mutations are computed.
            At the beginning it is composed of user-defined seeds, but during fuzzing this set can grow if a specific mutation discovers new edges not yet found.
        \item \textbf{Input Generator}: Composed of two main parts:
            \begin{itemize}
                \item \small \textbf{Seed Scheduling}: responsible for deciding which seed in the corpus to mutate.
                \item \small \textbf{Mutation Scheduling}: responsible for deciding which mutation or mutations to apply to the selected seed.
            \end{itemize}
        \item \textbf{Executor}: The component that actually executes the target binary.
            If the target binary has already been instrumented with Edge Coverage tracking instructions, it is executed as is;
            otherwise, in the case of \gls{DBI}, a dynamic translator must dynamically instrument the binary.
        \item \textbf{Coverage Collection}: In the case of Edge Coverage, this part is represented by the data structure used
            to keep track of which edges between \gls{BB}s have been taken and which have not.\\
            This information is used, in some tools such as AFL++ \cite{aflplusplus}, to guide Seed Scheduling,
            a technique that, given different inputs with the same level of coverage (i.e. number of edges visited),
            prioritizes the input that is smaller in size and takes less execution time.
        \item \textbf{Defect Monitor}: This component checks the exit status of the target to catch any crash or
            timeout of the target, or any other unexpected behavior.
            If none of these conditions are met, meaning that the target ends normally, the coverage is checked
            to see whether the current input discovered a new edge.
            If the input discovered a new edge, that input is added to the corpus.
        \end{enumerate}
\section{Binary Rewriting Background}\label{sec:biniary_rewriting}
    Binary rewriting consists in altering an already compiled executable without having access to the source code,
        while ensuring that the modified binary remains executable.
        The process of binary rewriting can be classified into Static Binary Rewriting, if the binary is stored in
        persistent memory while being rewritten, and Dynamic Binary Rewriting, if the binary is being executed.
        The process of binary rewriting can be divided into four main steps \cite{binaryrewriting}:
        \begin{enumerate}
            \item \textit{Parsing}: In this phase the main goal is to retrieve information about the binary.
                In particular, the output of this phase is a raw stream of instructions (in machine code) and data.
                This information is then used in the Analysis step.

            \item \textit{Analysis}: The information gained from the Parsing phase is analyzed to reconstruct the program structure.
            For example, the raw stream of machine code is passed through a disassembler that reconstructs the
            assembly instructions. A structural recovery is performed, which includes function detection and \gls{CFG} reconstruction,
            and labels and symbols are extracted from the raw stream of data.

            \item \textit{Transformation}: Once the \gls{CFG} has been reconstructed and function and data boundaries have been recovered,
                the binary rewriter can modify the binary at the \gls{IP}. These are user-defined locations in the binary where
                the rewriter must add, modify or delete an instruction. The \gls{IP} can be coarse- or fine-grained.
                In the first case, for example, the instrumentation could be inserted at branch level (at the beginning of each \gls{BB}).
                In the second case, instead, any instruction could potentially be modified.

            \item \textit{Code Generation}: In this last step, the changes made in the previous steps must be integrated into the
                binary while preserving its executability. This can be done in different ways. For example, the changes could be
                inserted into a new section of the executable, redirecting execution to the new code at the \gls{IP}.
                Alternatively, this step could also be performed by reassembling a completely new executable starting from the
                changes made at the \gls{IP}. This second approach is much harder, because it requires preserving all the
                references inside the executable and preserving the executable layout.
        \end{enumerate}
        \begin{figure}[ht]
            \caption{An example of a Static Binary Rewriting Workflow}
            \centering
            \includegraphics[width=200pt]{BinaryRewritingScheme.png}
            \label{fig:static_rewrite}
        \end{figure}
        In \cref{fig:static_rewrite}, a general scheme for static binary rewriting can be seen.
        In this case all the phases occur prior to execution, so the cost in terms of overhead
        is paid only once, before the execution. The main problems are disambiguating between code and data
        (i.e. deciding whether a stream of bytes represents machine code instructions or constants placed inline in the .text section)
        and resolving indirect branches \cite{binaryrewriting,andriesse2016}.
        This problem is generally undecidable and can only be solved using heuristics.
        There are three main families of static transformation:
        \begin{itemize}
            \item \textbf{Direct}: Consists in directly modifying the instruction. Works well with RISC systems with instructions having
                the same length. Generally fails when adding or removing instructions, because this does not preserve the offset into the .text section.
            \item \textbf{Minimal-invasive}: The new code is inserted into a new section of the binary, and at the \gls{IP}
                an unconditional jump to the new section is placed.
                The main limitation of this approach arises when the new section is far from the \gls{IP} and a near jump (necessary to reach the new section)
                uses a different number of bytes than the instruction at the instrumentation point.
                In fact, in x86-64 assembly a short jump uses 2 bytes while a near jump uses 5 bytes.
                So, if a short jump is needed at the instrumentation point for alignment reasons,
                but the new section is too far away and a near jump is needed instead, it is necessary to add a \textbf{trampoline}.
                This approach can be seen in \cref{fig:trampoline_instrumentation}.
                Trampoline-based instrumentation is one of the most common approaches to binary rewriting \cite{binaryrewriting}.

            \item \textbf{Full-translation}: The binary is lifted into an \gls{IR} that makes it possible to add an \gls{IP} at
                instruction level. The \gls{IR} is then reassembled into machine code after the changes are made.
        \end{itemize}
        \begin{figure}[ht]
            \caption{An example of Static Instrumentation using the trampoline technique}
            \centering
            \includegraphics[width=400pt]{trampolineInstrumentation.png}
            \label{fig:trampoline_instrumentation}
        \end{figure}
        With regard to Dynamic Binary Rewriting, the Analysis and Transformation phases are executed at runtime for each instruction.
        In this case, none of the problems that arise in Static Binary Rewriting occur, since the binary is being executed,
        so indirect branches can be resolved using the values in the registers, and there is no need to
        disambiguate between constants inside the .text section and instructions.
        The main drawback of this approach is the prohibitive overhead caused by the dynamic analysis and transformation of each instruction,
        together with the fact that only the taken branches can be modified.
        In fact, if a branch is not actually taken during execution, its instructions never reach the rewriter.\\
        \\
        A specific type of binary rewriting is Binary Instrumentation.
        Binary instrumentation uses binary rewriting techniques to add instructions that track binary behaviour,
        log events, or, in the case of this thesis, insert edge coverage tracing instructions.
\section{ELF Loading and Relocations}\label{sec:loading_relocations}
    \gls{ELF} is the standard executable format under Linux operating systems.
The work done in this thesis heavily involves the \gls{ELF} structure, with particular focus on some sections 
used for relocation, alongside the process of executing an \gls{ELF}.
\subsection{ELF Structure}
    In this subsection the process of \gls{ELF} execution is discussed, with a focus on the parts of the execution flow 
    that will influence this work.\\
    The \textbf{ELF Header} is a structure located at the beginning of the \gls{ELF} and contains some important information 
    about the type of \gls{ELF}, the target operating system \dots.
    Alongside this information there is the location of two important data structures inside the \gls{ELF}:
    \begin{itemize}
        \item The \textbf{Program Header Table}, which defines the \textit{Execution View}
        \item The \textbf{Section Header Table}, which defines the \textit{Linking View}
    \end{itemize}
    The \textbf{Program Header Table} contains information needed to create a runtime image of the program in memory, via 
    the \textbf{Segments} \textit{(i.e. which parts of the file must be loaded into memory, which shared 
    objects are needed \dots)}.
    Each segment contains the access flags for that part of memory (i.e. Read/Write/Execute).
    The \textbf{Section Header Table} contains information about the \textbf{Sections} that compose the binary, 
    together with their mapping to the segments (e.g. which segment the \textit{.text} section is mapped to).
    Among all the sections, some contain specific structures used to perform specific tasks:
    \begin{itemize}
        \item \textit{\textbf{.dynamic}: The dynamic section contains information used by the Dynamic Linker 
            (e.g. the dependencies, the interpreter location \dots)}.
        \item \textit{\textbf{.dynsym}: List of symbols defined and used by the binary. 
            If the binary is not stripped, there is also the \textbf{.symtab} section, which contains a more complete list of symbols}.
        \item \textit{\textbf{.hash or .gnu.hash}: Structures used for efficient symbol lookup during symbol resolution}\footnote{
            The \textbf{.gnu.hash} section is the more efficient version of \textbf{.hash}.
        }.
        \item \textit{\textbf{.rela.dyn and .rela.plt}: Both are relocation tables. The first is the standard one, 
        the second is used in the case of \textbf{Lazy Binding}}.
        \item \textit{\textbf{.shstrtab, .dynstr and .strtab}: Hold lists of strings referenced by other parts of the \gls{ELF}.
            The first holds section names, the second holds the names of the symbols contained in the \textbf{.dynsym} section, 
            and the last one holds the names of the entire \textbf{.symtab} section}.
        \item \textit{\textbf{.got and .got.plt}: Contain the \gls{GOT}, filled by the dynamic linker with the 
            addresses of external functions and variables. The first is the standard one, the second is used by the \gls{PLT}}.
        \item \textit{\textbf{.plt}: Contains the \gls{PLT}, which is an executable section made up of trampolines. In particular, it 
            contains stubs for functions that might be located in another \gls{ELF} or \gls{SO}. 
            The stub references the corresponding \gls{GOT} entry containing the resolved address of the function.
            This section is used in the case of \textbf{Lazy Binding}}.
        \item \textit{\textbf{.init and .init\_array}: Sections used for initialization. \textbf{.init} contains initialization code, 
            while \textbf{.init\_array} contains the addresses of initialization functions}.
        \item \textit{\textbf{.fini and .fini\_array}: Sections used for termination handling. 
            \textbf{.fini} contains termination code, 
            while \textbf{.fini\_array} contains the addresses of termination functions}.
    \end{itemize}
\subsection{ELF Linking}
    Before actually digging deep into \gls{ELF} loading, the linking process must be reviewed.
    The \textbf{Linking phase} consists of joining together different pieces of code belonging to different files.
    For example, if a program uses the \texttt{printf} function, it needs to take the code of printf from the libc library.
    This process of joining together the executable that uses printf with the printf code located in another file 
    is called \textbf{Linking}. Linking can be of two different types: 
    \begin{itemize}
        \item \textbf{Static linking}: The needed code (e.g. the instructions of printf) is copied inside the binary 
        at \textit{Link Time}. This causes the size of the binary to be bigger (because the code 
        of external functions actually lies inside it) and consequently wastes RAM space (if 10 binaries all use the 
        printf function, the code of printf is copied inside all 10 binaries).
        \item \textbf{Dynamic linking}: The needed code is not copied inside the binary; instead, the binary contains information (i.e. 
        the address of the external function or variable) used to locate the needed resource.
        The resolution happens at runtime, when the Dynamic Linker finds and loads the library into the process's virtual address space and 
        resolves the references to the symbols. This allows different processes to share the same \gls{SO}.
    \end{itemize}
    Since nowadays the most commonly used approach is \textbf{Dynamic Linking}, the following discussion will focus on it.\\
    As stated before, the linking process can be used for both external variables and external functions.
    In both cases the process is the same, since both functions and variables can be exported as symbols.\\
    When dealing with functions there is the possibility of using \textbf{Lazy Binding}.
    This technique consists of resolving the symbol only during the first invocation of the function, rather than at load time 
    as happens with variables.
    In order to achieve \textbf{Lazy Binding} a stub section is used (i.e. \textbf{.plt}).
    When a program uses a library function, a \gls{PLT} entry is added to the .plt section.
    \begin{figure}[ht]
    \centering
            \begin{subfigure}[h]{0.7\textwidth}
            \centering
            \includegraphics[width=\linewidth]{LazyBindingExample.png}
            \caption{An example of puts function resolution using Dynamic Linking and Lazy Binding before the actual resolution}
            \label{fig:lazy_resolution_plt_1}
        \end{subfigure}
        \vspace{0.5cm}
        \begin{subfigure}[h]{0.7\textwidth}
            \centering
            \includegraphics[width=\linewidth]{LazyBindingExampleAfterResolution.png}
            \caption{Same binary as \cref{fig:lazy_resolution_plt_1} but with symbol resolved}
            \label{fig:lazy_resolution_plt_2}
        \end{subfigure}
        \caption{An example of Lazy Binding involving .plt and .got}
        \label{fig:lazy_resolution}
    \end{figure}
    As can be seen in \cref{fig:lazy_resolution}, the relocation is resolved lazily upon the first invocation of the function. 
    When control is transferred to the corresponding \gls{PLT} entry, the entry performs an indirect jump through 
    its associated \gls{GOT} entry. Before the symbol is resolved, this \gls{GOT} entry contains an address that 
    redirects execution back to the \gls{PLT}, eventually reaching the \gls{PLT} resolver stub. 
    The resolver then invokes the dynamic linker, which performs the symbol lookup, resolves the relocation, 
    and updates the \gls{GOT} entry with the resolved address of the function, as can be seen in \cref{fig:lazy_resolution_plt_2}. 
    Subsequent calls can therefore jump directly to the resolved address through the \gls{GOT}.

\subsection{ELF Loading}
    The next useful step in order to understand the chosen approach consists of going into detail on the ELF Loading phase.
    There are 3 main actors that allow a binary to be executed properly:
    \begin{itemize}
        \item \textbf{Kernel}: responsible for the actual \texttt{execve} invocation and 
            allocation of the different memory segments. Moreover, it is the one that executes 
            the dynamic linker used to resolve dependencies prior to the execution of the binary.
        \item \textbf{Dynamic Linker}: Loads all the needed libraries into the process memory, 
            resolves the external dependencies, prepares the .got and .plt, and executes the initializers of the libraries.
        \item \textbf{Target Binary}: Invoked by the Dynamic Linker, it prepares the environment for the main function and, 
            in the end, runs it.
    \end{itemize}

    \begin{figure}[ht]
        \centering
        \includegraphics[width=300pt]{dynamicallyLinkedProgramLoad.png}
        \caption{Dynamically linked ELF loading procedure}
        \label{fig:dynamicallyLinkedProgramLoad}
    \end{figure}
    All these procedures can be seen in detail in \cref{fig:dynamicallyLinkedProgramLoad}.
    The dynamic linker part is the one on which I will focus in this dissertation.
    The dynamic linker:
    \begin{itemize}
        \item \textit{Loads all the DT\_NEEDED entries}: these are the dependencies (i.e. the \gls{SO}s needed by the binary).
            For example, if the binary uses an external function contained in \textit{libc.so}, the library from which the 
            function is exported must be added as a dependency.
        \item \textit{Resolves the relocations}: a focus on relocations will be presented in the next subsection.
        \item \textit{Arranges the \gls{GOT} and \gls{PLT}}.
        \item \textit{Executes all the initializer functions (constructors) for each dependency}: For each \textit{DT\_NEEDED} 
            entry found by the dynamic linker, it executes all the functions contained in the \textit{.init\_array} section of 
            that entry. These are functions declared with \texttt{\_\_attribute\_\_((constructor))}.
        \item \textit{Finds and executes the binary entry point}.
    \end{itemize}

\subsection{ELF Symbols and Relocations}
    The dynamic linker uses symbols to link different ELFs together.
    \begin{itemize}
        \item A component (an \gls{ELF} or a \gls{SO}) can export a symbol representing a variable or a function inside it.
        \item Another component can reference the exported symbol and use it (if the first component is a dependency of the second).
        \item The linker acts as the intermediary between the two components, responsible for matching the symbol names 
            (the exported one and the referenced one) and storing the address of the exported symbol in the second component's address space 
            in such a way that this component can access it.
    \end{itemize}
    All the symbols belonging to an \gls{ELF} (both imported and exported) are stored in the \textit{.dynsym} section.
    \begin{figure}[ht]
        \centering
        \includegraphics[width=\linewidth]{dynsym_example.png}
        \caption{Example of a \textit{.dynsym} table for a shared object}
        \label{fig:dynsym_example}
    \end{figure}
    \begin{figure}[ht]
        \centering
        \includegraphics[width=\linewidth]{reladynexample.png}
        \caption{Example of a \textit{.rela.dyn} relocation table extracted from a binary}
        \label{fig:reladyn_example}
    \end{figure}
    In \cref{fig:dynsym_example} a list of symbols extracted from a \gls{SO} can be seen.
    The first 4 symbols are external symbols that must be resolved. This can be seen because the Value column is 0 and Ndx is UND (undefined).
    The last 7, instead, are symbols declared inside the shared object.
    When the visibility is set to \textit{DEFAULT}, this means that the symbol can be visible outside the shared object.

    The process of resolving symbol references is called \textbf{Relocation}. The relocation algorithm is shown in 
    \cref{alg:relocation_resolution}.
    Relocating a symbol means resolving the symbol (i.e. finding the address of that symbol in the original object), and writing 
    that address inside the binary that has the relocation.
    Specifically, the symbol address is written inside the \gls{GOT}.
    In \cref{fig:reladyn_example}, the address at which the symbol address must be written can be seen in the ``Offset'' column.
    \begin{algorithm}
        \caption{Relocation resolution at ELF load time}
        \label{alg:relocation_resolution}
        \begin{algorithmic}[1]
            \REQUIRE $L = [obj_1,\dots,obj_n]$ \COMMENT{The objects that have a dependency (DT\_NEEDED) must already be loaded in memory}
            \ENSURE for each $obj \in L$, the respective GOT entry contains the resolved runtime address of the external symbol.
            
            \FOR{$obj \in L$}
                \STATE $Reloc\_table \gets obj.\texttt{GET\_SECTION(".rela.dyn")}$
                \FOR{$reloc \in Reloc\_table$}
                    \IF{$reloc.\text{type} = \texttt{R\_*\_RELATIVE}$}
                        \STATE $addr \gets obj.\texttt{loadBias} + r.\text{addend}$ 
                        \COMMENT{\textcolor{green!60!black}{The symbol does not need to be resolved. The address must be adjusted in order to enable ASLR}}
                    \ELSE
                        \IF{$reloc.\text{type} \in \{\texttt{R\_*\_GLOB\_DAT} , \texttt{R\_*\_JUMP\_SLOT}\}$} 
                            \STATE $sIdx \gets reloc.\text{symbol\_index}$
                            \STATE $symbolName \gets obj.\texttt{dynstr}[\, obj.\texttt{dynsym}[sIdx].\text{name} \,]$ 
                            \COMMENT{\textcolor{green!60!black}{The name of the symbol is read from \textit{.dynsym}}}
                            \STATE $target \gets \textsc{ResolveSymbol}(symbolName, L, obj)$
                            \IF{$target = \text{NULL}$}
                                \STATE \textbf{error}: undefined symbol $symbolName$ 
                            \ELSE
                                \STATE $(defObj, defIdx) \gets target$ \COMMENT{\small \textcolor{green!60!black}{\textit{defObj} is the object that exports the symbol, \textit{defIdx} is the index in the .dynsym table of \textit{defObj}}}
                                \STATE $addr \gets defObj.\texttt{loadBias} + defObj.\texttt{dynsym}[defIdx].\text{value}$
                                \COMMENT{\small \textcolor{green!60!black}{The address of the resolved symbol is computed as base\_address (of the object that exports the symbol) plus the offset at which the symbol is stored}}
                            \ENDIF
                        \ENDIF
                    \ENDIF
                    \STATE write $addr$ in $obj.\texttt{got}$ at offset $reloc.\text{offset}$ \COMMENT{\small \textcolor{green!60!black}{The destination address at which to write the resolved symbol address is specified in the relocation entry}}
                \ENDFOR
            \ENDFOR
        \end{algorithmic}
        
    \end{algorithm}
    \begin{algorithm}
        \caption{\textsc{ResolveSymbol} function used by \cref{alg:relocation_resolution}}
        \label{alg:symbol_resolution}
        \begin{algorithmic}[1]
            \REQUIRE name $name$, object list $L$, object requesting the resolution $requester$
            \ENSURE pair $(obj, idx)$ of the first symbol found, or NULL if the symbol has not been defined by any object.
            \FOR{$obj \in L$ \COMMENT{\small \textcolor{green!60!black}{The order is: first the executable, then all the DT\_NEEDED dependencies in order of definition}}}
                \STATE $idx \gets \textsc{GnuHashLookup}(name,\ obj.\texttt{gnuHash},\ obj.\texttt{dynsym},\ obj.\texttt{dynstr})$
                \COMMENT{\small \textcolor{green!60!black}{As seen before, the hashes are used in order to speed up the symbol lookup}}
                \IF{$idx \neq \text{NOT\_FOUND}$ \textbf{\&\&} $obj.\texttt{dynsym}[idx].\text{st\_shndx} \neq \texttt{SHN\_UNDEF}$}
                    \STATE \textbf{return} $(obj, idx)$ 
                \ENDIF
            \ENDFOR
            \STATE \textbf{return} NULL
        \end{algorithmic}
    \end{algorithm}
\section{Frame pointer omission and Red Zone}\label{sec:preamble}
    
The classic function preamble, inserted by the compiler at the beginning of a function, can be seen in \cref{lst:classic_preamble}.
\begin{lstlisting}[caption={Classic function preamble supposing a function that uses 0x30 bytes of stack},label=lst:classic_preamble]
push rbp 
mov rbp, rsp
sub rsp, 0x30
\end{lstlisting}
Every function has its own \textbf{stack frame}, that is the stack zone where all the local informations used by the function 
lies. Specifically it contains local variables, parameters passed to the function, \dots
The stack frame is delimited by the value of two registers:
\begin{itemize}
    \item \textbf{Base Pointer} : point to the base of the activation frame (high addresses)
    \item \textbf{Stack Pointer} : point to the top of the stack (low addresses)
\end{itemize}
When a function is invoked, the stack frame of the new function must be allocated. This is done by the \cref{lst:classic_preamble}.
The effect on the stack could be seen in \cref{fig:preamble1}. If the function had a local variable, that variable 
would be referenced with respect to the base pointer (e.g. \texttt{[rbp - 0x8]}).
\begin{figure}[ht]
    \centering
    \includegraphics[width=\linewidth]{preamble1.png}
    \caption{Example of what happend during a function call}
    \label{fig:preamble1}
\end{figure}
    The base pointer in this case is called \textbf{frame pointer} and it is used as a stable reference to the stack frame.
Infact, during the execution of a function, the stack pointer could change 
(a \texttt{push} or a \texttt{pop} instruction occur), and this case the stack pointer would not be a valid reference to the 
variables anymore.\\
If the compiler knows a priori that the stack pointer will not change during function execution, could optimize the function 
preamble with a technique called \textbf{frame pointer omission}.
This technique consist in not use the base pointer as reference to the local variables, but using directly the stack 
pointer. This allow the function preamble to be shrink to the one shown in \cref{lst:fp_omission_preamble}
\begin{lstlisting}[caption={Function preamble in case of frame pointer omission supposing a function that uses 0x30 bytes of stack},label=lst:fp_omission_preamble]
sub rsp, 0x30
\end{lstlisting}
In this case the access to local variable is expressed with respect to the stack pointer 
(e.g. \texttt{[rsp + 0x8]}), as can be seen in \cref{fig:preamble2}.\\
\begin{figure}[ht]
    \centering
    \includegraphics[width=0.55\linewidth]{preamble2.png}
    \caption{Example of what happend during a function call in case of frame pointer omission}
    \label{fig:preamble2}
\end{figure}
\begin{center}
    \textbf{red zone}
\end{center}
The red zone is composed of 128 bytes above the rsp that are reserved by the ABI System V AMD64.
The user code that is compliant with the ABI could use this reserved memory to save temporary data without actually 
modify the stack pointer and allocate a stack frame.
Since the stack grows toward low addresses, this zone is overwritten by the stack frame of any function that is called from 
the actual one, so the red zone could be used especially for leaf functions (functions that does not call other functions).\\
The compiler could choose to bypass entirely the function premable and access the local variable with respect to 
the stack pointer (e.g.\texttt{[rsp -0x8]}) as shown in \cref{fig:preamble3}
\begin{figure}[ht]
    \centering
    \includegraphics[width=0.25\linewidth]{preamble3.png}
    \caption{Example of function call that does not allocate space of the stack bu access the variable saved on the red zone}
    \label{fig:preamble3}
\end{figure}
\section{Problem Statement}\label{sec:problem_statement}
    Coverage-Guided Grey-box fuzzing is one of the most effective techniques for automatic vulnerability discovery in software
    \cite{zafl,Boehme2016},
    due to its ability to generate a wide range of inputs based on the feedback obtained through Code Coverage information,
    which guides the fuzzer towards unexplored regions of code and therefore helps it find hidden vulnerabilities.
    In fact, \textit{Black-box} fuzzing lacks a precise mutation strategy based on execution monitoring, while
    \textit{White-box} fuzzing incurs high overheads due to Symbolic or Concolic execution.
    Furthermore, the latter approach also requires the availability of the source code in the general case, 
    which is not always available.
    So \textit{Grey-box} fuzzing seems to be the most promising trade-off between precision of input mutation and overhead.
    However, \textit{Grey-box} fuzzing also has its own drawbacks.
    As a matter of fact, the effectiveness of this approach strongly depends on the availability of the source code.
    In fact, both \gls{SBI} and \gls{DBI} have many limitations.
    The main limitation of \gls{DBI} is the huge amount of overhead (higher than 600\% with respect to the compile-time 
    instrumentation in the case of AFL++ in QEMU mode \cite{zafl}), which makes the exploration of a substantial part
    of the valid input space almost impossible.
    Instead, the main difficulty of \gls{SBI} arises in modifying already compiled binaries.
    There are several problems regarding Binary Rewriting, such as Indirect Control Transfers (e.g. \texttt{call rax}),
    shifts in binary sections, or instructions that require rewriting the instructions that reference these sections, and many others,
    as seen in the previous \hyperref[sec:biniary_rewriting]{Section 1.2}.\\
    For all these reasons, the preferred technique nowadays is \textit{Compile-time Edge Coverage instrumentation} with fuzzers such as
    AFL++, LibAFL \cite{libafl}, or libFuzzer \cite{libfuzzer}.
    However, in practice, for reasons of industrial secrecy (i.e. proprietary software) or due to precompiled libraries, applications are
    often closed source, and so Compile-time Edge Coverage instrumentation becomes unusable.\\
    In these scenarios, binary-only fuzzing is the only viable approach.
    Although \gls{SBI} presents many challenges, it nonetheless proves to be the most fruitful area,
    given its low overheads \cite{zafl}.
    In this thesis I will address the problem
    of statically instrumenting a binary with Edge Coverage instructions in a way that makes the instrumented binary compatible with AFL++,
    which has been chosen as the fuzzer due to its popularity and because it represents the current state of the art, implementing
    most of the recently discovered techniques demonstrated by Andrea Fioraldi et al. \cite{aflplusplus}.\\
    The instrumentation must add as little overhead as possible with respect to compile-time overheads and must
    preserve the correctness of the binary. This is important in fuzzing, because lower overhead means more executions per second,
    so a larger portion of the valid input space can be explored.
\section{State of the Art}\label{sec:state_of_the_art}
    In this section I will present some of the current state-of-the-art solutions in binary-only grey-box fuzzing,
    distinguishing between \gls{SBI} and \gls{DBI} approaches.

    \begin{itemize}
        \item \underline{\textbf{\gls{DBI}}}
        \begin{itemize}
            \item \textbf{AFL++ in QEMU/FRIDA mode} \cite{aflplusplus}: AFL++ in QEMU mode uses QEMU in \textit{User-mode Emulation},
                which consists in translating guest machine code (i.e. guest \gls{BB}s) into an intermediate representation
                called \gls{TB} and subsequently appending the coverage tracing instructions at the end of it.
                This approach is based on emulating the entire CPU using QEMU's virtualized environment. As stated before,
                AFL in QEMU mode has overheads on the order of 600\% with respect to compile-time instrumentation \cite{zafl}.\\
                AFL++ in FRIDA mode \cite{aflplusplus-frida-mode}, instead, is based on shared library injection, in which a Frida \cite{frida} component
                called \textbf{Stalker} intercepts any \gls{BB} that is about to be executed, copies it into a \textbf{Code Cache}
                and dynamically instruments the copy with coverage tracing instructions.
                Execution control is then transferred to the copy, which is executed and then re-intercepted by \textbf{Stalker} at the end.
                Unlike QEMU mode, it does not emulate the CPU, but instead executes within the memory space of the target program.
                AFL++ reports that FRIDA mode is approximately 2 to 5 times slower than compile-time instrumentation \cite{aflplusplus_binary_only_targets}.

                \item \textbf{DynamoRIO/Pin} \cite{DynamoRIOPaper,PinPaper}: Both these approaches share the same basic principle, called
                \textbf{\gls{JIT} code caching}. This consists in taking control of the thread when the target application starts,
                and from that moment on, every \gls{BB} is intercepted before it is executed and is copied into a private \textbf{Code Cache}.
                The instrumentation is then added to the copy, and control is transferred to the code cache.
                The instrumentation process is more or less identical to FRIDA's Stalker mechanism, but is not as optimized.
                Also in this case, the implemented version uses AFL \cite{zalewski2017afl} as the fuzzer \cite{heuse2018afldynamorio,heuse2018aflpin}.
        \end{itemize}
            \item \underline{\textbf{\gls{SBI}}}
            \begin{itemize}
                \item \textbf{Retrowrite} \cite{retrowrite}: This tool performs symbolic lifting on the code, replacing every reference
                    discovered during disassembly with a symbolic label. Doing this allows instructions to be added without
                    modifying the references, achieving reassembleable assembly.
                    While RetroWrite can add ASan instrumentation, it cannot natively add AFL-style edge coverage instructions.
                    To do so, the lifted assembly must be instrumented with the afl-gcc tool in order to produce the patched binary.
                    Moreover, community-reported issues \cite{retrowrite_issue28} document missing instrumentation points when
                    RetroWrite's symbolized output is compiled through afl-gcc, caused by a mismatch between RetroWrite's
                    label-naming convention and afl-as's basic-block detection heuristics.
                    Finally, RetroWrite only works with PIE, non-stripped binaries.
                \item \textbf{AFL-Dyninst} \cite{afldyninst}: AFL-Dyninst employs Dyninst's binary rewriting capabilities
                        to instrument an existing executable by inserting callbacks at selected basic blocks.
                        Dyninst's instrumentation mechanism is based on code patching and trampoline redirection \cite{Dyninst},
                        which causes many context switches and state saving/restoring operations.
                        Its performance is slightly better than \gls{DBI}, but this approach still reaches 500\% fuzzing overhead,
                        as demonstrated by Stefan Nagy et al. \cite{zafl}.
                
                \item \textbf{E9AFL} \cite{E9AFL} : E9AFL uses E9Patch \cite{E9Patch} to instrument binaries.
                    This is a framework that allow to modify existing isntruction without actually retrieve and 
                    analyze the whole \gls{CFG}. 
                    E9AFL modify existing instruction inserting trampolines to a code region that update the \gls{SHM}.
                    Since trampoline based instrumentation suffer of many overheads, this approach introduce many 
                    optimization techniques that reduce the overheads given by the trampolines. 
                    Alvin Charles et al. \cite{PeAR} shows that E9AFL has some limitations, especially regarding the 
                    capability of instrumenting correctly the target binaries.

                \item \textbf{ZAFL} \cite{zafl}: This tool uses two main phases: the first consists in
                    reconstructing an \gls{IR} using the Zipr binary rewriter.
                    The resulting \gls{IR} is then used by ZAX to perform an optimization on the \gls{IR}, find the instrumentation points,
                    and finally instrument these points. ZAFL works by inserting inline calls to a function that updates the shared map
                    at the beginning of each \gls{BB}.
                    Despite ZAFL's many improvements, Jifan Xiao et al. \cite{SystemCallPatternTracing} demonstrated that
                    it still causes crashes in 28\% of the binaries used in their cited paper.
            \end{itemize}
    \end{itemize}
    Before proceeding to the next section, in which I will discuss the main goals of this thesis, I want to discuss the core functioning of AFL and AFL++,
    given its importance for understanding my approach.
    \subsection{AFL Functioning}
    \cite{afl_technical_details}
        As previously stated, AFL and, in particular, AFL++ are among the most widely used grey-box
        fuzzing frameworks \cite{aflplusplus}.
        AFL++ represents the community-based continuation of the original AFL project.
        The core functioning is virtually the same and consists of a \gls{SHM} region shared between the fuzzer and the target.
        The \gls{SHM} size can vary, but the default is 65536 bytes.\\
        Each ``cell'' of the \gls{SHM} contains the number of times a given edge in the \gls{CFG} has been
            taken during the execution of the target. Since the \gls{SHM} has 65536 ``cells'', each ``cell''
            can contain a number from 0 to 255.
        The ``cell'' position is computed with a specific hash function applied to the IDs of the two \gls{BB}s
        that form the edge.
        The main Edge Tracing update rule is summarized in \cref{alg:afl_style_update}.
        \begin{algorithm}
            \caption{AFL-style SHM update rule}
            \label{alg:afl_style_update}
            \begin{algorithmic}[1]
                \STATE $idx \gets current\_bb\_id \oplus previous\_bb\_id$ 
                \STATE $shared\_map[idx] \gets shared\_map[idx]+1$ 
                \STATE $previous\_bb\_id \gets current\_bb\_id \gg 1$ 
            \end{algorithmic}
        \end{algorithm}
        Line 3 of \cref{alg:afl_style_update} is used to avoid updating the same \gls{SHM}
        location when an edge is traversed forward or backward.
        An example of this could be a simple \texttt{do-while} loop, as seen in \cref{alg:do-while} and \cref{fig:Do-While-CFG}.\\
        \begin{figure}[ht]
            \centering
            \begin{minipage}[h]{0.48\textwidth}
                \centering
                \begin{algorithm}[H]
                    \caption{Do-While Example}
                    \label{alg:do-while}
                    \begin{algorithmic}
                        \REPEAT
                            \STATE \COMMENT{Basic Block A}
                        \UNTIL{$Permanence\_Condition$}
                        \STATE \COMMENT{Basic Block C}
                    \end{algorithmic}
                \end{algorithm}
            Without 3 of \cref{alg:afl_style_update}, the hash of Edge1
            would result in $ID\_A \oplus ID\_B$, and the hash of Edge2 (if the permanence condition is not met) would result
            in $ID\_B \oplus ID\_A = ID\_A \oplus ID\_B$, due to the symmetry of XOR. But since the hash is the same, there would be a hash collision
            that would cause the same ``cell'' to be updated twice instead of once.
            \end{minipage}
            \hfill
            \begin{minipage}[h]{0.48\textwidth}
                \centering
                \includegraphics[width=100pt]{AFLUpdateAlg3InstructionWhy.png}
                \caption{Do-while CFG}
                \label{fig:Do-While-CFG}
            \end{minipage}
        \end{figure}
        The main idea behind AFL is to instrument each basic block of a target application with the shared map update rule.
        The fuzzer initializes the \gls{SHM} memory with zeros, runs the instrumented target, which updates the \gls{SHM}
        according to \cref{alg:afl_style_update} until it reaches the end, and then returns to the fuzzer.
        The \gls{SHM} ID is communicated to the target through an environment variable called \textit{\_\_AFL\_SHM\_ID}.\\
        At this point the fuzzer checks the \gls{SHM} it has just obtained, called \texttt{trace\_bits}, searching for any newly discovered edge.
        This is done by keeping a structure called \texttt{virgin\_bits}, which contains all the edges not yet discovered and
        is initialized with ones. The comparison is performed through $\texttt{trace\_bits} \land \texttt{virgin\_bits}$. If any position results in one after the logical AND,
        this means a new edge has been discovered. In this case, the \texttt{virgin\_bits} data structure is updated and
        the input just executed is saved in the corpus. Otherwise, a new input is chosen, mutated and then executed. \\
        This approach requires a \texttt{fork()} call followed by an \texttt{execve()} call for each input,
        causing significant overhead for each input execution.
        Two main optimizations have been implemented:
        \begin{itemize}
            \item \textbf{forkserver}
            \item \textbf{persistent mode}: \small (not shown in this thesis)
        \end{itemize}
        \vspace{0.5cm}
        \begin{center}
            \textbf{Forkserver}
        \end{center}
        The forkserver makes it possible to avoid calling \textbf{execve} for every input, performing it only once at the start of execution.
        The forkserver is instrumented before the main function of the target application.
        The AFL++ forkserver protocol can be summarized as follows:
        \begin{figure}[ht]
            \caption{AFL++ Forkserver Protocol}
            \centering
            \includegraphics[width=350pt]{AFLPlusPlusForkserverFunctioning.png}
            \label{fig:aflplusplusforkserver}
        \end{figure}
        \footnote{ \small \textit{The differences between the AFL forkserver protocol and the AFL++ forkserver protocol lie mainly in the messages exchanged between the forkserver and the fuzzer}}
        \begin{enumerate}
            \item The forkserver sends the ``Welcome'' message to the fuzzer (0x41464c01).
            \item The fuzzer replies with the status code. At this point the forkserver must check whether $status \oplus 0xffffffff = 0x41464c01$.
                    If it is equal, execution continues; otherwise, the forkserver stops here.
            \item The forkserver sends the execution options to the fuzzer.
                These can indicate that the target is running in persistent mode, or other optimizations.
            \item The forkserver sends 0x41464c01 to the fuzzer to end the handshake.
            \item The fuzzer writes a message to the forkserver (its content is discarded); the handshake is now finished.
                        This message is the ``go'' signal for the forkserver.
            \item The forkserver invokes \texttt{fork()}. The parent writes the PID of the child to the fuzzer and then invokes \texttt{waitpid()}.\\
                    The child continues execution of the main function until it ends.
            \item When the parent returns from the \texttt{waitpid()} function (i.e. the child has terminated), the exit status
            of the child is sent to the fuzzer, and the forkserver waits for another ``go'' from the fuzzer.
        \end{enumerate}
        The communication between the Fuzzer and the Target shown in \cref{fig:aflplusplusforkserver} occurs
        using two unidirectionally pipes that uses the file descriptor:
        \begin{itemize}
            \item \textit{FORKSRV\_FD (fd 198)}: used by the forkserver to read a message sent by the fuzzer.
            \item \textit{FORKSRV\_FD + 1 (fd 199)}: used by the forkserver to write a message to the fuzzer.
        \end{itemize}
        
        An approach worth mention in this section is the \textit{deferred forkserver} technique.
        It consist in execute the forkserver as later as possible in the target. 
        As shown previously, at each input, all the instructions after the forkserver must be executed by the child. 
        If the forkserver is invoked before some instructions that does not require tracing (for example the load of a 
        configuration file, or the initialization of an object), these instructions are executed for each input, 
        producing useless overheads.
\section{Goals and Challenges}\label{sec:goals-and-challenges}
    The main goal of this thesis is to develop a tool capable of inline, static edge-coverage instrumentation
that preserves binary correctness while achieving better performance with respect to \gls{DBI} and comparable to 
other state-of-the-art instrumentation technique.\\
I have chosen to implement the \gls{SBI} approach instead of \gls{DBI} due to its low overhead.
To do this, I will use the GRAPE binary rewriting framework, a novel approach for binary rewriting that preserves binary executability
while introducing small overheads compared to a trampoline-based approach.
This approach involves several challenges, mainly related to preserving the binary's correctness and behavior.
Moreover, the tool must be compatible with the already existing AFL++ fuzzer.
This means implementing the forkserver communication protocol inside the target binary without recompiling it.
This has created no small number of problems, mainly regarding the preservation of the ELF layout after adding new sections.
Compared to the ZAFL approach, the \gls{SHM} update instructions are placed directly inside the \gls{BB},
rather than in a function called at each \gls{BB}.
This approach drastically increases the size of the file, causing many instruction repetitions (since the \gls{SHM}
instructions are always the same), but was chosen because it does not cause a context switch at each \gls{BB}.
In fact, in order to call a function, the state of the registers must be saved and restored. Moreover, the stack must be allocated,
and the callee's RIP must be saved on the stack alongside the old base pointer of the stack.
Performing all these operations at each \gls{BB} would introduce significant overhead.